\SourceLeaf{121}%
\emph{Demonstr.} I. Sint $\alpha,\beta$ integri tales qui faciant $\alpha a+\beta b=1$, unde erit
\[
i=\alpha b-\beta a+m(\beta+\alpha i)
\]
Proposito itaque numero integro complexo $A+Bi$, habebimus
\[
A+Bi=A+(\alpha b-\beta a)B+m(\beta B+\alpha Bi)
\]
Quare denotando per $h$ residuum minimum positivum numeri $A+(\alpha b-\beta a)B$ secundum modulum $p$, statuendoque
\[
A+(\alpha b-\beta a)B=h+kp=h+m(ak-bki)
\]
erit
\[
A+Bi=h+m(\beta B+ak+(\alpha B-bk)i)
\]
sive
\[
A+Bi\equiv h(\mathrm{mod.}\ m).\qquad\text{Q. E. P.}
\]

II. Quoties eidem numero complexo duo numeri reales $h,h'$ secundum modulum $m$ congrui sunt, etiam inter se congrui erunt. Statuamus itaque $h-h'=m(c+di)$, unde fit
\[
(h-h')(a-bi)=p(c+di)
\]
adeoque
\[
(h-h')a=pc,\quad(h-h')b=-pd
\]
nec non, propter $a\alpha+b\beta=1$,
\[
h-h'=p(c\alpha-d\beta),\ \text{i. e.}\ h\equiv h'(\mathrm{mod.}\ p)
\]
Quapropter $h$ et $h'$, siquidem sunt inaequales, ambo simul in complexu numerorum $0,1,2,3\dotsfour p-1$ contenti esse nequeunt. Q. E. S.

\Article{41}
\textsc{Theorema.}\quad\emph{Secundum modulum complexum $m=.\,a+bi$, cuius norma $aa+bb=p$, et pro quo $a,b$ non sunt inter se primi, sed divisorem communem maximum $\lambda$ habent (quem positive acceptum supponimus), quilibet numerus complexus congruus est residuo $x+yi$ tali, ut $x$ sit aliquis numerorum $0,1,2,3\dotsfour\frac p\lambda-1$, atque $y$ aliquis horum $0,1,2,3\dotsfour\lambda-1$, et quidem unico tantum inter omnia $p$ residua, quae tali forma gaudent.}

\SourceLeaf{122}%
\emph{Demonstr.} I. Accipiendo integros $\alpha,\beta$ ita, ut fiat $\alpha a+\beta b=\lambda$, erit $\lambda i=\alpha b-\beta a+m(\beta+\alpha i)$. Iam sit $A+Bi$ numerus complexus propositus, $y$ residuum minimum positivum ipsius $B$ secundum modulum $\lambda$, atque $x$ residuum minimum positivum ipsius $A+(\alpha b-\beta a)\pdot\frac{B-y}{\lambda}$ secundum modulum $\frac p\lambda$, statuaturque
\[
A+(\alpha b-\beta a)\pdot\frac{B-y}{\lambda}=x+\frac p\lambda\pdot k
\]
Hinc erit
\[
\begin{aligned}
A+Bi-(x+yi)&=\frac p\lambda\pdot k+(B-y)i-(\alpha b-\beta a)\frac{B-y}{\lambda}\\
&=\frac p\lambda\pdot k+\frac{B-y}{\lambda}\pdot m(\beta+\alpha i)\\
&=\left(\frac a\lambda-\frac b\lambda\pdot i\right)km+\frac{B-y}{\lambda}(\beta+\alpha i)m
\end{aligned}
\]
i. e. per $m$ divisibilis, sive $A+Bi\equiv x+yi(\mathrm{mod.}\ m)$ Q. E. P.

II. Supponamus, secundum modulum $m$ eidem numero complexo congruos esse duos numeros $x+yi,x'+y'i$, qui proin etiam inter se congrui erunt secundum modulum $m$. A potiori itaque secundum modulum $\lambda$ congrui erunt, adeoque $y\equiv y'(\mathrm{mod.}\lambda)$. Quodsi igitur uterque $y,y'$ inter numeros $0,1,2,3\dotsfour\lambda-1$ contentus esse supponitur, necessario debet esse $y=y'$. Hoc pacto vero etiam fiet $x\equiv x'(\mathrm{mod.}\ m)$, i. e. $x-x'$ per $m$, adeoque $\frac{x-x'}\lambda$ integer per $\frac a\lambda+\frac b\lambda\pdot i$ divisibilis, sive
\[
\frac{x-x'}\lambda\equiv0\left(\mathrm{mod.}\frac a\lambda+\frac b\lambda\pdot i\right)
\]
Hinc autem, quum $\frac a\lambda,\frac b\lambda$ sint numeri inter se primi, concluditur per partem secundam theorematis art. praec., $\frac{x-x'}\lambda$ etiam per normam numeri $\frac a\lambda+\frac b\lambda\pdot i$, i. e. per numerum $\frac p{\lambda\lambda}$ divisibilem fore, adeoque $x-x'$ per $\frac p\lambda$. Quapropter si etiam uterque $x,x'$ in complexu numerorum $0,1,2,3\dotsfour\frac p\lambda-1$ contentus esse supponitur, necessario erit $x=x'$, sive residua $x+yi,x'+y'i$ identica. Q. E. S.

Ceterum sponte patet, huc quoque referendum esse casum, ubi modulus est numerus realis, puta $b=0$, et proin $\lambda=\pm a$, nec non eum, ubi modulus est numerus pure imaginarius, puta $a=0$, et proin $\lambda=\pm b$. In utroque casu habetur $\frac p\lambda=\lambda$.

\SourceLeaf{123}%
\Article{42}
Referendo itaque omnes numeros complexos secundum modulum datum inter se congruos ad eandem classem, incongruos ad diversas, omnino aderunt $p$ classes totum numerorum integrorum ambitum exhaurientes, denotante $p$ normam moduli. Complexus totidem numerorum e singulis classibus desumtorum exhibebit systema completum residuorum incongruorum, quale in artt. 40, 41 assignavimus. Et in hocce quidem systemate electio residuorum classes suas quasi repraesentantium innixa erat principio ei, ut in quavis classe adoptaretur residuum $x+yi$ tale, pro quo $y$ habeat valorem minimum, atque inter omnia, quibus idem valor minimus ipsius $y$ inest, id, pro quo valor ipsius $x$ est minimus, exclusis valoribus negativis tum pro $x$ tum pro $y$. Sed ad alia proposita aliis principiis uti conveniet, imprimisque notandus est modus is, ubi residua talia adoptantur, quae per modulum divisa offerunt quotientes simplicissimos. Manifesto si $\alpha+\beta i,\alpha'+\beta'i,\alpha''+\beta''i$ etc. sunt quotientes e divisione numerorum congruorum per modulum oriundi, differentiae tum quantitatum $\alpha,\alpha',\alpha''$ etc. inter se erunt numeri integri, tum differentiae inter quantitates $\beta,\beta',\beta''$ etc., patetque, semper adesse residuum unum, pro quo $\alpha$ et $\beta$ iaceant inter limites $0$ et $1$, limite priori incluso, posteriori excluso: tale residuum simpliciter vocamus residuum minimum. Si magis placet, loco illorum limitum etiam hi adoptari possunt $-\frac12$ et $+\frac12$ (altero admisso, altero excluso): residuum tali limitationi respondens \emph{absolute minimum} dicemus.

Circa\textperiodcentered haec residua minima offerunt se problemata sequentia.

\Article{43}
Residuum minimum numeri complexi dati $A+Bi$ secundum modulum $a+bi$, cuius norma $=p$, invenitur sequenti modo. Si $x+yi$ est residuum minimum quaesitum, erit $(x+yi)(a-bi)$ residuum minimum producti $(A+Bi)(a-bi)$ secundum modulum $(a+bi)(a-bi)$, i. e. secundum modulum $p$. Statuendo itaque
\[
aA+bB=Fp+f,\quad aB-bA=Gp+g
\]
ita ut $f,g$ sint residua minima numerorum $aA+bB,aB-bA$, secundum modulum $p$, erit
\SourceLeaf{124}%
\[
x+yi=\frac{f+gi}{a-bi}
\]
sive
\[
\begin{aligned}
x&=\frac{af-bg}{p}=A-aF+bG\\
y&=\frac{ag+bf}{p}=B-aG-bF
\end{aligned}
\]
Manifesto residua minima $f,g$ vel inter limites $0$ et $p-1$, vel inter hos $-\frac12p$ et $+\frac12p$ accipi debent, prout numeri complexi vel residuum simpliciter minimum vel absolute minimum desideratur.

\Article{44}
Constructio systematis completi residuorum minimorum pro modulo dato pluribus modis effici potest. Methodus prima ita procedit, ut primo determinentur limites, intra quos termini reales iacere debent, ac dein pro singulis valoribus intra hos limites sitis assignentur limites partium imaginariarum. Criterium generale residui minimi $x+yi$ pro modulo $a+bi$ in eo consistit, ut tum $ax+by=\xi$, tum $ay-bx=\eta$ iaceat inter limites $0$ et $aa+bb$, quoties de residuis simpliciter minimis agitur, vel inter limites $-\frac12(aa+bb)$ et $+\frac12(aa+bb)$, quoties residua absolute minima desiderantur, limite altero excluso. Regulae speciales distinctionem casuum, quos varietas signorum numerorum $a,b$ affert, requirerent, cui tamen evolvendae, quum nulli difficultati obnoxia sit, hic immorari supersedemus: sufficiat, methodi indolem per unicum exemplum exposuisse.

Pro modulo $5+2i$ residua simpliciter minima $x+yi$ ita comparata esse debent, ut tum $5x+2y=\xi$, tum $5y-2x=\eta$ aequetur alicui numerorum $0,1,2,3\dotsfour28$. Aequatio $29x=5\xi-2\eta$ ostendit, valores positivos ipsius $x$ maiores esse non posse quam $\frac{5\pdot28}{29}$, negativos abstrahendo a signo non maiores quam $\frac{2\pdot28}{29}$. Omnes itaque valores admissibiles ipsius $x$ erunt $-1,0,1,2,3,4$. Pro $x=-1$ debet esse $2y$ aequalis alicui numerorum $5,6,7\dotsfour33$, atque $5y$ alicui horum $-2,-1,0,1\dotsfour26$; hinc valor minimus ipsius $y$ est $+3$, maximus $+5$. Tractando perinde valores reliquos ipsius $x$, oritur sequens schema omnium residuorum minimorum:

\SourceLeaf{125}%
\[
\begin{array}{r|l}
x&y\\\hline
-1&3,4,5\\
0&0,1,2,3,4,5\\
+1&1,2,3\pdot4,5\pdot6\\
+2&1,2,3,4,5,6\\
+3&2,3,4\pdot5,6\\
+4&2,3,4
\end{array}
\]

Simili modo pro residuis absolute minimis, $\xi$ et $\eta$ alicui numerorum $-14,-13,-12\dotsfour+14$ aequales esse debent; hinc $29x$ nequit esse extra limites $-7\pdot14$ et $+7\pdot14$, adeoque $x$ alicui numerorum $-3,-2,-1,0,1,2,3$ aequalis esse debet. Pro $x=-3$ erit $2y=\xi-5x=\xi+15$ alicui numerorum $1,2,3\dotsfour29$ aequalis, $5y=\eta+2x=\eta-6$ autem alicui horum $-20,-19,-18\dotsfour+8$: hinc prodit pro $y$ valor unicus $+1$. Tractando eodem modo valores reliquos ipsius $x$, habemus schema omnium residuorum absolute minimorum:

\[
\begin{array}{r|l}
x&y\\\hline
-3&+1\\
-2&-2,-1,0,+1,+2\\
-1&-3,-2,-1,0,+1,+2\\
0&-2,-1,0,+1,+2\\
+1&-2,-1,0,+1,+2,+3\\
+2&-2,-1,0,+1,+2\\
+3&-1
\end{array}
\]

\Article{45}
In applicatione methodi secundae duos casus distinguere conveniet.

In casu priori, ubi $a$ et $b$ divisorem communem non habent, fiat $a\alpha+\beta b=1$, sitque $k$ residuum minimum positivum ipsius $\beta a-\alpha b$ secundum modulum $p$. Hinc aequationes identicae
\[
\begin{aligned}
a(\beta a-\alpha b)&=\beta p-b(\alpha a+\beta b),\\
b(\beta a-\alpha b)&=-\alpha p+a(\alpha a+\beta b)
\end{aligned}
\]
docent, esse $ak\equiv-b,bk\equiv a(\mathrm{mod.}\ p)$. Statuendo itaque ut supra $ax+by=\xi$,
